\chapter{总结与展望}

\section{工作总结}

本文面向精密舵机桌面测试场景，完成了一套基于分层架构的测试软件。舵机测试包含串口接入、执行对象确认、反馈采集、性能计算和结果复查等环节，单次指令发送只是其中的基础动作。论文的工作范围围绕这些已经实现的环节展开，未把缺少实测依据的 24 h 稳定性、完整端到端时延和图片导出流程写成已完成成果。

系统结构采用界面层、桥接层、协调层、后台任务层、串口通信层和测试核心分工。界面层保留操作和显示，串口、文件和后台任务由受控接口进入主进程，测试核心通过运行时接口读取反馈样本。该划分把高频反馈处理从界面交互中分离出来，也让后续章节可以分别讨论通信协议、使能检测、测试算法和结果保存。

通信与基础控制部分完成了两类协议帧的构造、解析和校验处理。连续串口字节流经过缓存和帧边界识别后，转化为上层可以直接读取的反馈样本。使能检测在正式测试前执行，用基线采样、小扰动激励和相关判别确认当前执行对象的响应状态。

测试功能按实际操作顺序组织。通信性能测试先检查收发链路；控制性能预测试估计噪声标准差、保守调节时间和平均运动速度；时域测试依据这些参数生成位置路径并计算响应指标；频域测试利用扫频激励和相关法估计幅值、相位和伯德图指标。参数传递、任务进度、结果保存和历史加载贯穿这四类测试。

工程实现和验证部分说明了软件的当前可用范围。自动化测试检查协议解析、测试核心、信号发生器和状态切换等逻辑；实机检查在当前联调条件下覆盖手动收发、高频自动发送和使能检测。据此，本文可以确认软件已经具备当前阶段舵机测试、结果显示和结构化留档能力。

\section{系统设计与实现特点}

从实现结果看，本系统的一个特点是把测试过程拆成可交接的几段。串口接入和协议解析先给出反馈样本，使能检测再确认执行对象状态，四类测试在此基础上读取同一类反馈数据，最后由保存模块处理结果文件。每一段都有明确输入和输出，界面组件不需要直接处理底层串口字节。

通信层的处理方式也为后续扩展留下了余地。同步字、帧长、字段偏移、序号和校验规则集中在通信模块内，协议差异不会散落到时域测试、频域测试或数据显示组件中。若后续增加新的协议版本，主要修改位置应集中在帧构造、帧解析和反馈样本映射部分。

参数传递是测试流程中的另一条线索。通信测试给出链路状态，预测试给出噪声、调节时间和速度估计，时域与频域测试再用这些参数决定采样规模、等待时间和路径长度。用户仍可以修改参数，后端也会对非法值进行校验。正式测试的参数来源因而更容易追溯。

结果管理没有停留在界面显示。通信统计、预测试参数、时域曲线、频域曲线和指标摘要会进入共享结果状态，并可以保存为结构化文件。历史文件加载后能够恢复图表和指标显示，这为后续复查、对比和异常定位提供了数据基础。

\section{不足与展望}

当前工作仍受到实验条件和工程周期限制。系统已经完成当前联调条件下的自动化测试和实机检查，但还没有长时间连续运行报告。串口链路、后台任务、界面刷新和结果保存都可能在更长测试中暴露新的问题。后续需要先确定测试工况和判定标准，再记录异常恢复、资源占用和数据完整性。

实时性评价也需要单独补充。本文没有把完整端到端时延列为已验证指标。预测试和时域测试中的调节时间反映的是舵机控制响应，不能替代从界面操作、控制发送、串口反馈、数据解析到界面刷新这一整条链路的时延测量。后续应明确起点、终点和采样方法，再分别统计各段耗时。

结果管理仍有完善空间。当前系统已经具备结构化结果保存和历史加载能力，图片导出、长期批量归档和自动报告生成还没有形成完整流程。后续可围绕测试批次、设备编号、协议版本、测试参数和结果摘要整理归档规则，使多次测试之间的比较更加方便。

协议适配和测试覆盖范围也需要继续扩展。当前设计已经把协议差异集中到通信层，把反馈数据整理为统一样本。后续仍应在更多舵机型号、通信配置和测试工况下检查该设计是否适用。时域和频域测试中的判据、参数推荐和异常提示，也需要结合更多实机数据逐步修正。

本文完成了精密舵机测试软件的分层设计、通信处理、四类测试和结果管理实现，并在当前条件下完成了阶段性验证。下一步工作应集中在长期稳定性、端到端时延、数据归档和协议扩展上。
